% ============================================================
\chapter{Nonlinear Systems --- Linearization and Center Manifold}
\label{ch:nonlinear}
% ============================================================

We now leave the reassuring world of linear systems and enter the infinitely richer realm of nonlinear dynamics. Here, surprises abound: a pendulum can oscillate or spin, a population can explode or collapse, an electrical circuit can generate self-sustained oscillations. The Hartman-Grobman theorem (1960) assures us that near a \emph{hyperbolic} equilibrium, the nonlinear phase portrait qualitatively resembles the linearised one. But what happens when eigenvalues fall on the imaginary axis? This is where \emph{centre manifold theory}, developed by Victor Pliss (1964) and Jack Carr (1981), isolates the essential dynamics on a submanifold of minimal dimension.

\section{Introduction}

Nonlinear systems exhibit a phenomenological richness far beyond linear
systems: multiple equilibria, limit cycles, chaos. To understand
\emph{local} behavior near an equilibrium, we use linearization
(Hartman-Grobman theorem) and, in degenerate cases, center manifold theory.

\section{The Hartman-Grobman Theorem}

\begin{theorem}[Hartman-Grobman]
\label{thm:hartman-grobman}
Let $\dot{x} = f(x)$ with $f \in C^1$, $f(0) = 0$, and $A = Df(0)$.
If $A$ is \textbf{hyperbolic} (no eigenvalues on the imaginary axis), then
there exists a homeomorphism $h$ defined in a neighborhood of $0$ that
conjugates the flow of $\dot{x} = f(x)$ to the linear flow $\dot{y} = Ay$:
\[
    h(\varphi_t(x)) = e^{At}\, h(x)
\]
for $t$ sufficiently small and $x$ in a neighborhood of $0$.
\end{theorem}

\begin{remark}
The Hartman-Grobman theorem guarantees that the local phase portrait topology
is entirely determined by the linearization when the equilibrium is
hyperbolic. The homeomorphism $h$ is generally not differentiable.
\end{remark}

\begin{definition}[Hyperbolic equilibrium]
An equilibrium $x^*$ is \textbf{hyperbolic} if the Jacobian matrix
$Df(x^*)$ has no eigenvalue with zero real part.
\end{definition}

\begin{definition}[Non-hyperbolic equilibrium]
An equilibrium is \textbf{non-hyperbolic} (or \textbf{degenerate}) if at
least one eigenvalue of $Df(x^*)$ has zero real part. In this case,
linearization does not determine stability: additional tools are needed.
\end{definition}

\section{Local Invariant Manifolds}

\begin{theorem}[Stable and unstable manifolds]
\label{thm:stable-manifolds}
Let $x^* = 0$ be a hyperbolic equilibrium of $\dot{x} = f(x)$ with
$A = Df(0)$. Let $E^s$ and $E^u$ be the stable and unstable subspaces of
$A$. Then there exist local manifolds:
\begin{itemize}
    \item $W^s_{\mathrm{loc}}(0)$: \textbf{local stable manifold}, tangent
    to $E^s$ at $0$, of dimension $\dim E^s$;
    \item $W^u_{\mathrm{loc}}(0)$: \textbf{local unstable manifold}, tangent
    to $E^u$ at $0$, of dimension $\dim E^u$.
\end{itemize}
These manifolds are of class $C^k$ if $f \in C^k$ and are locally invariant
under the flow.
\end{theorem}

\begin{definition}[Global manifolds]
The \textbf{global stable and unstable manifolds} are defined by:
\[
    W^s(0) = \{x \in U : \varphi(t, x) \to 0 \text{ as } t \to +\infty\},
\]
\[
    W^u(0) = \{x \in U : \varphi(t, x) \to 0 \text{ as } t \to -\infty\}.
\]
\end{definition}

\begin{example}[Saddle in dimension 2]
For the system $\dot{x} = -x + x^2$, $\dot{y} = y$, the origin is a saddle
with $E^s = \{y = 0\}$, $E^u = \{x = 0\}$. The local stable manifold is the
curve $y = 0$, $\abs{x} < 1$, and the unstable manifold is the $x = 0$ axis.
\end{example}

\section{Center Manifold Theory}

When the equilibrium is non-hyperbolic, the essential dynamics is captured
by the \textbf{center manifold}.

\begin{theorem}[Existence of the center manifold]
\label{thm:center-manifold}
Let $\dot{x} = f(x)$ with $f(0) = 0$ and $A = Df(0)$. Suppose $A$ has $n_s$
eigenvalues with negative real part, $n_u$ with positive real part, and $n_c$
with zero real part, where $n_s + n_u + n_c = n$. Then there exists a local
manifold $W^c_{\mathrm{loc}}(0)$:
\begin{itemize}
    \item of dimension $n_c$;
    \item tangent to $E^c$ at $0$;
    \item locally invariant under the flow;
    \item of class $C^k$ if $f \in C^k$ (but not necessarily analytic even
    if $f$ is analytic).
\end{itemize}
\end{theorem}

\begin{warning}[title=\textbf{Non-uniqueness of the center manifold}]
Unlike stable and unstable manifolds, the center manifold is generally
\textbf{not unique}. Two different center manifolds may agree to all orders
but differ by exponentially small terms.
\end{warning}

\begin{theorem}[Reduction to the center manifold]
\label{thm:center-reduction}
If $n_u = 0$ (no unstable direction), the stability of the equilibrium
$x^* = 0$ for the full system is determined by the stability of $0$ in the
dynamics restricted to the center manifold $W^c$.
\end{theorem}

\subsection{Practical Computation of the Center Manifold}

Consider the system in canonical form:
\[
    \dot{x} = Cx + F(x, y), \qquad \dot{y} = Sy + G(x, y),
\]
where $x \in \R^{n_c}$, $y \in \R^{n_s}$, $C$ has all eigenvalues on the
imaginary axis, $S$ has all eigenvalues with negative real part, and $F$, $G$
are at least second order.

The center manifold is the graph of a function $y = h(x)$ with $h(0) = 0$,
$Dh(0) = 0$. The function $h$ satisfies the partial differential equation:
\[
    Dh(x)\,[Cx + F(x, h(x))] = Sh(x) + G(x, h(x)).
\]

In practice, one expands $h$ in a Taylor series and identifies coefficients
order by order.

\begin{example}[Center manifold in dimension 2]
Consider the system
\[
    \dot{x} = xy, \qquad \dot{y} = -y + x^2.
\]
The origin is a non-hyperbolic equilibrium:
$A = \begin{pmatrix} 0 & 0 \\ 0 & -1 \end{pmatrix}$, $E^c = \{y = 0\}$,
$E^s = \{x = 0\}$.

We seek $y = h(x) = ax^2 + bx^3 + \cdots$ Substituting into the center
manifold equation:
\[
    h'(x) \cdot x\,h(x) = -h(x) + x^2.
\]
At order 2: $0 = -ax^2 + x^2$, so $a = 1$.
At order 3: the left side contributes $h'(x) \cdot x \cdot h(x) = 2x \cdot x \cdot x^2 + \cdots = 2x^4 + \cdots$,
which is order 4. So $b = 0$ at order 3.

The dynamics on the center manifold is:
\[
    \dot{x} = x \cdot h(x) = x \cdot x^2 + \cdots = x^3 + \cdots
\]
Since $\dot{x} = x^3 > 0$ for $x > 0$ and $\dot{x} < 0$ for $x < 0$,
the equilibrium is \textbf{unstable}.
\end{example}

\section{Normal Forms}

\begin{definition}[Poincar\'e normal form]
A \textbf{normal form} is a simplification of the nonlinear system obtained
through successive polynomial changes of variables that eliminate
non-resonant terms.
\end{definition}

\begin{theorem}[Normal form theorem]
Let $\dot{x} = Ax + f_2(x) + f_3(x) + \cdots$ with $f_k$ homogeneous of
degree $k$. For each $k \geq 2$, there exists a polynomial change of
variable $x = y + \phi_k(y)$ that eliminates the non-resonant terms of
order $k$ in $f_k$.
\end{theorem}

\begin{definition}[Resonance]
A monomial $x^m = x_1^{m_1} \cdots x_n^{m_n}$ in the $j$-th component is
\textbf{resonant} if
\[
    \langle m, \lambda \rangle - \lambda_j = 0,
\]
where $\lambda = (\lambda_1, \ldots, \lambda_n)$ is the eigenvalue vector
and $\langle m, \lambda \rangle = m_1\lambda_1 + \cdots + m_n\lambda_n$.
\end{definition}

\section{Smooth Linearization}

\begin{theorem}[Sternberg]
If $x^* = 0$ is hyperbolic and the eigenvalues of $Df(0)$ satisfy no
resonance relation, then there exists a local $C^\infty$ diffeomorphism
conjugating the nonlinear system to the linear system $\dot{y} = Ay$.
\end{theorem}

\begin{remark}
Sternberg's theorem improves upon Hartman-Grobman by providing a
diffeomorphism (not merely a homeomorphism) under non-resonance conditions.
\end{remark}

\section{Application: Population Dynamics}

\begin{example}[Competition between two species]
The Lotka-Volterra competition model is
\[
    \dot{x}_1 = x_1(1 - x_1 - \alpha x_2), \qquad
    \dot{x}_2 = x_2(1 - x_2 - \beta x_1),
\]
with $\alpha, \beta > 0$. The equilibrium points are $(0,0)$, $(1,0)$,
$(0,1)$ and, if $\alpha \neq 1$ and $\beta \neq 1$, the coexistence point
\[
    x^* = \left(\frac{1-\alpha}{1-\alpha\beta}, \frac{1-\beta}{1-\alpha\beta}\right).
\]
Linearization at each equilibrium determines the local dynamics.
\end{example}

\begin{codeblock}[title=\textbf{Stable and unstable manifolds of a saddle}]
\begin{minted}{python}
import numpy as np
from scipy.integrate import solve_ivp
import matplotlib.pyplot as plt

def system(t, state):
    x, y = state
    return [x - x**2, -y + x*y]

fig, ax = plt.subplots(figsize=(8, 6))

# Stable manifold (integrate backward from near saddle)
for eps in [0.01, -0.01]:
    sol = solve_ivp(lambda t, s: [-v for v in system(t, s)],
                    [0, 5], [0, eps], max_step=0.01)
    ax.plot(sol.y[0], sol.y[1], 'b-', lw=2,
            label='$W^s$' if eps > 0 else '')

# Unstable manifold (integrate forward)
for eps in [0.01, -0.01]:
    sol = solve_ivp(system, [0, 5], [eps, 0], max_step=0.01)
    ax.plot(sol.y[0], sol.y[1], 'r-', lw=2,
            label='$W^u$' if eps > 0 else '')

ax.set_xlim(-1, 2); ax.set_ylim(-1, 2)
ax.set_xlabel('$x$'); ax.set_ylabel('$y$')
ax.legend(); ax.grid(True, alpha=0.3)
ax.set_title("Stable and unstable manifolds")
plt.tight_layout()
plt.savefig("stable_unstable_manifolds.pdf")
plt.show()
\end{minted}
\end{codeblock}

\section{Exercises}

\begin{exercise}[Hartman-Grobman]
Verify that the equilibrium $(0,0)$ of the system
$\dot{x} = -x + y^2$, $\dot{y} = -2y + x^2$ is hyperbolic.
What is the type of the local phase portrait?
\end{exercise}

\begin{exercise}[Center manifold]
Compute the center manifold and determine the stability of the origin for
\[
    \dot{x} = -x^2y, \qquad \dot{y} = -y + x^2.
\]
\end{exercise}

\begin{exercise}[Transcritical bifurcation]
Study the system $\dot{x} = \mu x - x^2$, $\dot{y} = -y$ near
$(\mu, x, y) = (0, 0, 0)$ by center manifold reduction.
\end{exercise}

\begin{exercise}[Normal form]
Put the system $\dot{x} = -y + ax^2 + bxy$, $\dot{y} = x + cx^2 + dxy$
into normal form up to order 3.
\end{exercise}

\begin{exercise}[Global manifolds]
Numerically plot the stable and unstable manifolds of the equilibrium
$(1, 0)$ for $\dot{x} = x(1-x) - xy$, $\dot{y} = y(x - 1)$.
\end{exercise}

\begin{exercise}[Non-uniqueness]
Show that for the system $\dot{x} = x^2$, $\dot{y} = -y$, the center
manifold is not unique by exhibiting two functions $h_1(x)$ and $h_2(x)$
satisfying the center manifold equation.
\end{exercise}
